%!TEX program = lualatex
% Arquivo: relatorio_pratica_5_2algarismos.tex

\documentclass[a4paper,12pt]{article}

% =========================================================
% Idioma e tipografia (LuaLaTeX)
% =========================================================
\usepackage[brazilian]{babel}
\usepackage{fontspec}
\setmainfont{Latin Modern Roman}
\usepackage{microtype}

% =========================================================
% Layout
% =========================================================
\usepackage[top=3cm,bottom=2.5cm,left=3cm,right=2.5cm]{geometry}
\usepackage{setspace}
\onehalfspacing
\usepackage{indentfirst}
\setlength{\parindent}{1.25cm}

% =========================================================
% Matemática, unidades e fórmulas químicas
% =========================================================
\usepackage{amsmath,amssymb}
\usepackage{siunitx}
\DeclareSIUnit{\ppm}{ppm}
\sisetup{
  locale = DE,
  output-decimal-marker = {,},
  per-mode = symbol,
  exponent-product = \cdot
}
\usepackage[version=4]{mhchem}

% =========================================================
% Tabelas, listas, figuras e desenhos
% =========================================================
\usepackage{booktabs}
\usepackage{array}
\usepackage{enumitem}
\setlist[itemize]{noitemsep,topsep=3pt}
\setlist[enumerate]{noitemsep,topsep=3pt}
\usepackage{graphicx}
\usepackage{float}
\graphicspath{{./}{./imagens/}{./figuras/}}

\usepackage{xcolor}
\usepackage{tikz}
\usetikzlibrary{arrows.meta,calc,positioning,shapes.geometric}
\usepackage[most]{tcolorbox}

% =========================================================
% Links
% =========================================================
\usepackage{hyperref}
\usepackage{xurl}
\hypersetup{
  colorlinks=true,
  linkcolor=black,
  urlcolor=blue,
  citecolor=black,
  pdfauthor={Israel Carneiro; José Mauro Xavier Elsas; Miguel Gayani},
  pdftitle={Prática 5 -- Determinação de cloreto em água do mar}
}

% =========================================================
% Cores e comandos
% =========================================================
\definecolor{quadroazul}{RGB}{40,91,140}
\definecolor{quadrovermelho}{RGB}{177,69,60}
\definecolor{quadrocinza}{RGB}{245,245,245}

\newcommand{\subtitulo}[1]{%
  \par\medskip
  {\color{quadroazul}\bfseries\large #1}\par\smallskip
}

% =========================================================
% Dados da capa
% =========================================================
\newcommand{\Instituicao}{Universidade do Estado do Rio de Janeiro -- UERJ}
\newcommand{\Curso}{Química VIII / Oceanografia}
\newcommand{\Professor}{Professor: Ítalo Vieira}
\newcommand{\Aluno}{Alunos: Israel Carneiro\\
José Mauro Xavier Elsas\\
Miguel Gayani}
\newcommand{\Cidade}{Rio de Janeiro}
\newcommand{\Ano}{2026}

\begin{document}

% =========================================================
% Capa
% =========================================================
\begin{titlepage}
\centering
{\large \Instituicao \par}
\vspace{0.5cm}
{\large \Curso \par}
\vspace{1.0cm}

{\LARGE \textbf{Prática 5} \par}
\vspace{0.6cm}
{\Large \textbf{DETERMINAÇÃO DE CLORETO\\ EM ÁGUA DO MAR }}

\vspace{1cm}

\IfFileExists{uerj.jpg}{%
  \begin{figure}[H]
    \centering
    \includegraphics[width=0.3\linewidth]{uerj.jpg}
  \end{figure}
}{%
  \begin{tcolorbox}[
    width=0.32\textwidth,
    colback=white,
    colframe=black!35,
    boxrule=0.5pt,
    arc=0mm
  ]
  \centering
  \vspace{1.2cm}
  \textit{Logotipo da UERJ}
  \vspace{1.2cm}
  \end{tcolorbox}
}

\vfill

\begin{flushright}
\Aluno\\
\Professor
\end{flushright}

\vspace{1.0cm}
{\Cidade{} -- \Ano \par}
\end{titlepage}

% =========================================================
\section{Introdução}

Os cloretos são íons originados de diferentes sais dissolvidos nas águas e seu teor varia de acordo com o tipo de água, tendo valores limites estabelecidos por órgãos reguladores como o CONAMA. Diferentes tipos de métodos são empregados na análise de cloretos incluindo os métodos clássicos, como os métodos volumétricos de precipitação. Entre os métodos volumétricos de precipitação, os mais comuns são os que empregam solução padrão de nitrato de prata e são amplamente utilizados na determinação de haletos
(cloreto, brometo e iodeto), tiocianato (SCN-) e cianeto (CN-), com formação de sais de prata pouco solúveis.

Baseado nos diferentes tipos de indicadores, existem três métodos: o Método de Mohr, o de Volhard e o de Fajans. O método de Mohr será empregado na presente aula prática.


% =========================================================
\section{Objetivos}

\begin{itemize}
  \item padronizar a solução de Nitrato de Prata;
  \item executar a titulação de Cloreto em água do mar.

\end{itemize}

% =========================================================
\section{Material utilizado}

\begin{enumerate}[label=\alph*.]
  \item béquer de \qty{400}{\milli\litre};
  \item balão volumétrico de \qty{250}{\milli\litre};
  \item pipeta volumétrica \qty{25}{\milli\litre};
  \item pipeta volumétrica \qty{10}{\milli\litre};
  \item bureta de \qty{50.00}{\milli\litre};
  \item suporte universal e garra para bureta;
  \item frascos erlenmeyer de \qty{250}{\milli\litre};
 \end{enumerate}

% =========================================================
\section{Reagentes e soluções}

\begin{enumerate}[label=\alph*.]
  \item solução de \ce{NaCl}, \qty{0.0141}{\mole\per\litre};
  \item solução de \ce{AgNO3}  \qty{0.014}{\mole\per\litre};
  \item solução indicadora de Cromato de Potássio a \qty{0.002}{\mole\per\litre};
  \end{enumerate}

% =========================================================
\section{Procedimento experimental}

\subsection{Padronização da solução de \ce{AgNO3}}

A bureta foi rinsada e preenchida com a solução de \ce{AgNO3}, de
concentração nominal aproximada de \qty{0.014}{\mole\per\litre}. Para
cada determinação, transferiu-se uma alíquota de
\qty{10.00}{\milli\litre} da solução padrão de \ce{NaCl}, de concentração
\qty{0.0141}{\mole\per\litre}, para um erlenmeyer. Em seguida,
adicionou-se \qty{1}{\milli\litre} da solução indicadora de cromato de
potássio.

A solução de \ce{AgNO3} foi adicionada lentamente, sob agitação. Próximo
ao ponto final, as paredes internas do erlenmeyer foram lavadas com água
destilada e o titulante passou a ser adicionado gota a gota. Considerou-se
como ponto final o aparecimento persistente da coloração
vermelho-alaranjada, decorrente da formação de \ce{Ag2CrO4}. O
procedimento foi realizado em triplicata.

\begin{figure}[H]
\centering
\begin{tikzpicture}[>=Latex,thick,font=\small]
  \draw[quadroazul] (-0.15,1.8) rectangle (0.15,5.5);
  \draw[quadroazul] (-0.30,5.5) -- (0.30,5.5);
  \draw[quadroazul] (-0.15,1.8) -- (0,1.3) -- (0.15,1.8);
  \draw[quadroazul] (0,1.3) -- (0,0.9);
  \node[quadroazul,right=5mm] at (0.15,4.7)
    {\ce{AgNO3} $\approx$ \qty{0.014}{\mole\per\litre}};
  \fill[quadroazul!50] (0,0.65) circle (0.07);
  \draw[quadrovermelho]
    (-0.35,0.35) .. controls (-1.25,-0.15) and (-1.15,-1.05) .. (-0.55,-1.35)
    -- (0.55,-1.35)
    .. controls (1.15,-1.05) and (1.25,-0.15) .. (0.35,0.35)
    -- (0.35,1.10) -- (-0.35,1.10) -- cycle;
  \fill[yellow!20]
    (-0.70,-0.95) .. controls (-0.90,-0.65) and (-0.85,-0.35) .. (-0.55,-0.20)
    -- (0.55,-0.20)
    .. controls (0.85,-0.35) and (0.90,-0.65) .. (0.70,-0.95) -- cycle;
  \node[align=center,below=5mm] at (0,-1.35)
    {\qty{10.00}{\milli\litre} de \ce{NaCl}\\
     + \qty{1}{\milli\litre} de indicador};
  \draw[->,quadrovermelho] (2.2,0.6) -- (0.55,-0.1);
  \node[quadrovermelho,align=left,right] at (2.25,0.6)
    {amarelo $\longrightarrow$\\vermelho-alaranjado};
\end{tikzpicture}
\caption{Esquema da padronização da solução de \ce{AgNO3}.}
\label{fig:padronizacao_agno3}
\end{figure}

A reação de precipitação ocorre na proporção estequiométrica $1:1$:

\begin{equation}
  \ce{Ag+ (aq) + Cl- (aq) -> AgCl (s)}.
  \label{eq:precipitacao_agcl}
\end{equation}

Depois que o cloreto é consumido, o primeiro excesso de íons prata reage
com o cromato:

\begin{equation}
  \ce{2 Ag+ (aq) + CrO4^2- (aq) -> Ag2CrO4 (s)}.
  \label{eq:indicador_cromato}
\end{equation}

\subsection{Titulação de cloreto na amostra de água do mar}

Pipetaram-se \qty{5.00}{\milli\litre} da amostra de água do mar para um
balão volumétrico de \qty{100.00}{\milli\litre}. O volume foi completado
com água destilada até o traço de referência e a solução foi
homogeneizada.

Uma alíquota de \qty{10.00}{\milli\litre} da solução diluída foi
transferida para um erlenmeyer e recebeu \qty{1}{\milli\litre} da solução
indicadora de cromato de potássio. A mistura foi titulada com a solução
padronizada de \ce{AgNO3} até o aparecimento persistente da coloração
vermelho-alaranjada. A determinação foi realizada em triplicata.

\begin{figure}[H]
\centering
\begin{tikzpicture}[>=Latex,thick,font=\small]
  \draw[quadroazul] (-4.5,1.0) -- (-4.5,2.2) -- (-4.8,2.55)
    -- (-4.2,2.55) -- (-4.5,2.2);
  \draw[quadroazul] (-4.5,1.0)
    .. controls (-5.6,0.5) and (-5.5,-1.2) .. (-4.5,-1.4);
  \draw[quadroazul] (-4.5,1.0)
    .. controls (-3.4,0.5) and (-3.5,-1.2) .. (-4.5,-1.4);
  \node[align=center] at (-4.5,-2.0)
    {\qty{5.00}{\milli\litre} de água do mar\\
     diluídos para \qty{100.00}{\milli\litre}};
  \draw[->,quadroazul] (-2.7,0) -- (-0.6,0)
    node[midway,above,align=center]{pipetar\\\qty{10.00}{\milli\litre}};
  \draw[quadrovermelho] (0.0,1.0) -- (0.0,0.45) -- (-0.9,-1.15)
    -- (1.0,-1.15) -- (0.1,0.45) -- (0.1,1.0);
  \draw[quadrovermelho] (0.0,1.0) -- (0.1,1.0);
  \node[align=left,anchor=west] at (1.45,0)
    {$+\ \qty{1}{\milli\litre}$ de indicador\\
     titular com \ce{AgNO3} padronizado};
  \node[quadrovermelho] at (0.75,-2.0)
    {amarelo $\longrightarrow$ vermelho-alaranjado};
\end{tikzpicture}
\caption{Diluição da amostra e titulação de uma alíquota.}
\label{fig:titulacao_amostra}
\end{figure}

O ensaio em branco não foi realizado. Por essa razão, os cálculos a
seguir empregam diretamente os volumes lidos na bureta, sem correção de
branco.

% =========================================================
\section{Resultados e cálculos}

\subsection{Padronização da solução de \ce{AgNO3}}

Os volumes consumidos na padronização estão apresentados na
Tabela~\ref{tab:padronizacao}.

\begin{table}[H]
  \centering
  \caption{Volumes de \ce{AgNO3} consumidos na padronização.}
  \label{tab:padronizacao}
  \begin{tabular}{c S[table-format=2.2]}
    \toprule
    \textbf{Determinação}
      & {\textbf{Volume de \ce{AgNO3}} (\si{\milli\litre})} \\
    \midrule
    1 & 10,4 \\
    2 & 10,3 \\
    3 & 10,2 \\
    \midrule
    \textbf{Média} & 10,3 \\
    \bottomrule
  \end{tabular}
\end{table}

O volume médio de titulante foi

\begin{equation}
  \overline{V}_{\ce{AgNO3}}
  =\frac{10.4+10.3+10.2}{3}
  =\boxed{\qty{10.30}{\milli\litre}}.
  \label{eq:media_padronizacao}
\end{equation}

O desvio-padrão amostral foi de \qty{0.10}{\milli\litre} e o
desvio-padrão relativo foi de \qty{0.97}{\percent}.

No ponto de equivalência,

\begin{equation}
  n_{\ce{Ag+}}=n_{\ce{Cl-}},
\end{equation}

portanto,

\begin{equation}
  C_{\ce{AgNO3}}V_{\ce{AgNO3}}
  =C_{\ce{NaCl}}V_{\ce{NaCl}}.
  \label{eq:padronizacao}
\end{equation}

A concentração padronizada da solução de \ce{AgNO3} foi

\begin{align}
  C_{\ce{AgNO3}}
  &=
  \frac{C_{\ce{NaCl}}V_{\ce{NaCl}}}
       {\overline{V}_{\ce{AgNO3}}} \nonumber\\
  &=
  \frac{
    \left(\qty{0.0141}{\mole\per\litre}\right)
    \left(\qty{10.00}{\milli\litre}\right)}
    {\qty{10.30}{\milli\litre}} \nonumber\\
  &=\boxed{\qty{0.014}{\mole\per\litre}}.
  \label{eq:concentracao_agno3}
\end{align}

\subsection{Titulação da amostra diluída}

Nos cálculos subsequentes foi mantido internamente o valor não
arredondado da concentração de \ce{AgNO3}. Apenas os resultados
apresentados foram arredondados para dois algarismos significativos.

Os volumes consumidos nas três titulações da amostra estão apresentados
na Tabela~\ref{tab:titulacao_amostra}.

\begin{table}[H]
  \centering
  \caption{Volumes de \ce{AgNO3} consumidos na titulação da amostra diluída.}
  \label{tab:titulacao_amostra}
  \begin{tabular}{c S[table-format=2.2]}
    \toprule
    \textbf{Determinação}
      & {\textbf{Volume de \ce{AgNO3}} (\si{\milli\litre})} \\
    \midrule
    1 & 20,6 \\
    2 & 20,2 \\
    3 & 20,3 \\
    \midrule
    \textbf{Média} & 20,37 \\
    \bottomrule
  \end{tabular}
\end{table}

O volume médio foi

\begin{equation}
  \overline{V}_{\ce{AgNO3},\text{amostra}}
  =\frac{20.6+20.2+20.3}{3}
  =\boxed{\qty{20.37}{\milli\litre}}.
  \label{eq:media_amostra}
\end{equation}

O desvio-padrão amostral foi de \qty{0.21}{\milli\litre} e o
desvio-padrão relativo foi de \qty{1.02}{\percent}.

\subsection{Concentração de cloreto em \texorpdfstring{\si{\mole\per\litre}}{mol/L}}

Na alíquota titulada, a proporção entre \ce{Ag+} e \ce{Cl-} é $1:1$.
Assim, a concentração de cloreto na solução diluída foi

\begin{align}
  C_{\ce{Cl-},\text{diluída}}
  &=
  \frac{
    C_{\ce{AgNO3}}\,
    \overline{V}_{\ce{AgNO3},\text{amostra}}}
    {V_{\text{alíquota}}} \nonumber\\
  &=
  \left[
    \frac{
      \left(\num{0.0141}\right)
      \left(\num{10.00}\right)}
      {\num{10.30}}
  \right]
  \frac{\num{20.37}}{\num{10.00}}
  \,\si{\mole\per\litre} \nonumber\\
  &=\boxed{\qty{0.028}{\mole\per\litre}}.
  \label{eq:cloreto_diluida}
\end{align}

Como \qty{5.00}{\milli\litre} da amostra original foram diluídos para
\qty{100.00}{\milli\litre}, o fator de diluição foi

\begin{equation}
  F_D=\frac{100.00}{5.00}=20.
\end{equation}

Logo, a concentração de cloreto na amostra original foi

\begin{align}
  C_{\ce{Cl-},\text{amostra}}
  &=F_D\,C_{\ce{Cl-},\text{diluída}} \nonumber\\
  &=20\,C_{\ce{Cl-},\text{diluída}} \nonumber\\
  &=\boxed{\qty{0.56}{\mole\per\litre}}.
  \label{eq:cloreto_molar}
\end{align}

\subsection{Teor de cloreto em porcentagem massa/volume}

Adotando a massa molar do cloreto como
$M_{\ce{Cl-}}=\qty{35.45}{\gram\per\mole}$, obteve-se

\begin{align}
  \rho_{\ce{Cl-}}
  &=
  C_{\ce{Cl-},\text{amostra}}M_{\ce{Cl-}} \nonumber\\
  &=
  C_{\ce{Cl-},\text{amostra}}
  \left(\qty{35.45}{\gram\per\mole}\right) \nonumber\\
  &=\boxed{\qty{2.0e1}{\gram\per\litre}}.
\end{align}

Esse resultado corresponde a \qty{2.0}{\gram} de cloreto em
\qty{100}{\milli\litre} de amostra. Portanto,

\begin{equation}
  \boxed{\text{Teor de \ce{Cl-}}=\qty{2.0}{\percent}\ (m/V)}.
  \label{eq:cloreto_percentual}
\end{equation}

\subsection{Teor de cloreto em ppm}

Para soluções aquosas, tomando a aproximação
$\qty{1}{\milli\gram\per\litre}\approx\qty{1}{\ppm}$,

\begin{align}
  [\ce{Cl-}]
  &=
  \rho_{\ce{Cl-}}
  \left(\frac{\qty{1000}{\milli\gram}}{\qty{1}{\gram}}\right)
  \nonumber\\
  &=\boxed{\qty{2.0e4}{\milli\gram\per\litre}}
  \approx\boxed{\qty{2.0e4}{\ppm}}.
  \label{eq:cloreto_ppm}
\end{align}

% =========================================================
\section{Discussão}

A padronização forneceu concentração de
\qty{0.014}{\mole\per\litre} para a solução de \ce{AgNO3}, valor
próximo à concentração nominal de \qty{0.014}{\mole\per\litre}. O
desvio-padrão relativo de \qty{0.97}{\percent} indica boa proximidade
entre as três determinações.

Na titulação da amostra diluída, o volume médio de \ce{AgNO3} foi de
\qty{20.37}{\milli\litre}, com desvio-padrão relativo de
\qty{1.02}{\percent}. Após a correção pelo fator de diluição 20, a
concentração de cloreto na água do mar foi estimada em
\qty{0.56}{\mole\per\litre}, equivalente a
\qty{2.0e1}{\gram\per\litre}, \qty{2.0}{\percent} $(m/V)$ ou
aproximadamente \qty{2.0e4}{\ppm}.

O método de Mohr baseia-se na menor solubilidade do \ce{AgCl} em relação
ao \ce{Ag2CrO4}. Enquanto houver cloreto disponível, o íon prata
adicionado precipita predominantemente como \ce{AgCl}. Após o consumo do
cloreto, um pequeno excesso de \ce{Ag+} reage com o cromato, formando o
precipitado vermelho-alaranjado que sinaliza o ponto final.

A principal limitação dos resultados apresentados é o fato de o ensaio
em branco não ter sido realizado. A apostila prevê essa correção porque
uma pequena quantidade adicional de \ce{AgNO3} é necessária para produzir
\ce{Ag2CrO4} em quantidade visualmente perceptível. Sem a subtração do
branco, o volume atribuído ao cloreto pode ficar ligeiramente
superestimado e, consequentemente, também o teor calculado. Outras fontes
de incerteza incluem a identificação visual do ponto final, a leitura da
bureta e eventuais perdas ou erros durante a diluição e a transferência
das alíquotas.

% =========================================================
\section{Conclusão}

A solução de \ce{AgNO3} foi padronizada com solução de
\ce{NaCl} a \qty{0.0141}{\mole\per\litre}. A partir do volume médio de
\qty{10.30}{\milli\litre}, obteve-se a concentração de
\qty{0.014}{\mole\per\litre} para o titulante.

Na análise da amostra, o consumo médio foi de
\qty{20.37}{\milli\litre} de \ce{AgNO3}. Considerando a diluição de
\qty{5.00}{\milli\litre} para \qty{100.00}{\milli\litre}, o teor de
cloreto na água do mar foi calculado como

\begin{equation*}
  \boxed{
  C_{\ce{Cl-}}=\qty{0.56}{\mole\per\litre}
  =\qty{2.0}{\percent}\ (m/V)
  \approx\qty{2.0e4}{\ppm}
  }.
\end{equation*}

Os resultados das três determinações apresentaram boa repetibilidade.
Como o ensaio em branco não foi realizado, os valores foram calculados
diretamente a partir dos volumes experimentais medidos, sem a correção
prevista no procedimento da apostila.

% =========================================================
\section*{Referência}

LAGO, Dalva Cristina Baptista do; GOIS, Jefferson Santos de; LANGONE,
Marta Antunes Pereira; SENNA, Lilian Ferreira de.
\textit{Química VIII -- Oceanografia: aulas práticas}.
Rio de Janeiro: Universidade do Estado do Rio de Janeiro,
Instituto de Química, Departamento de Química Analítica, 2017.

\end{document}
